\chapter{Cross-Cloud Service Equivalence}
\index{cross-cloud equivalence}\index{service mapping}\index{multi-cloud}%
This appendix is a living reference that maps conceptually equivalent services and
ideas across the seven platforms studied in this book series: Amazon Web Services
(AWS), Microsoft Azure, Microsoft Fabric, Google Cloud Platform (GCP), MongoDB,
Snowflake, and Databricks. The names differ, but the underlying data-engineering
concepts---durable object storage, tiered lifecycle economics, immutability, managed
relational engines, elastic compute, open table formats, and disciplined data
organization---recur everywhere.

\begin{notebox}
This appendix grows with each chapter. Every new chapter contributes the equivalence
tables introduced in its cross-cloud callout, so by the end of the book you will hold
a single consolidated Rosetta Stone for moving fluently between clouds. Where a clean
one-to-one mapping does not exist, the prose notes the closest analog and the caveat.
\end{notebox}

\section{Object and Data-Lake Storage}
\index{object storage!equivalence}\index{storage classes!equivalence}%
The foundational layer covered in Chapter~1. Every platform offers durable, HTTP-addressable
object storage with tiered pricing, lifecycle transitions, versioning, and write-once-read-many
(WORM) controls.

\begin{center}
\footnotesize
\begin{tabular}{@{}p{2.8cm}p{3.2cm}p{3.2cm}p{3.2cm}@{}}
\toprule
\textbf{Capability} & \textbf{AWS} & \textbf{Azure} & \textbf{GCP} \\
\midrule
Object storage & Amazon S3 & ADLS Gen2 / Blob & Cloud Storage \\
Hot tier & S3 Standard & Hot & Standard \\
Infrequent access & S3 Standard-IA & Cool & Nearline \\
Archive (instant) & Glacier Instant Retrieval & Cold & Coldline \\
Deep archive & Glacier Deep Archive & Archive & Archive \\
Lifecycle rules & Lifecycle policies & Lifecycle management & Object Lifecycle Mgmt \\
Versioning & S3 Versioning & Blob versioning & Object Versioning \\
Immutability (WORM) & Object Lock & Immutable blob storage & Bucket Lock \\
Cross-region copy & CRR / SRR & Object replication & Dual/multi-region + Transfer \\
Encryption at rest & SSE-S3 / SSE-KMS & SSE + CMK (Key Vault) & Google-managed / CMEK \\
\bottomrule
\end{tabular}
\end{center}

\noindent The remaining platforms approach storage differently.
\textbf{Microsoft Fabric} exposes a single logical lake, \emph{OneLake}, and uses
\emph{shortcuts} to virtualize S3, ADLS Gen2, and Google Cloud Storage in place without copying.
\textbf{Snowflake} abstracts storage entirely: table data lives in cloud object storage as
immutable, columnar \emph{micro-partitions} that the service manages on your behalf.
\textbf{MongoDB} persists documents through the \emph{WiredTiger} storage engine, balancing an
internal cache against the operating-system page cache rather than exposing storage tiers.
\textbf{Databricks} sells no object storage at all: it layers Delta Lake tables and Unity Catalog
\emph{volumes} over the same S3, ADLS Gen2, or GCS buckets you already own, so the storage
conversation is identical to the AWS one---plus a table format on top.

\section{Managed Relational Databases}
\index{relational databases!equivalence}%
The focus of Chapter~2 for the three hyperscalers. Each provides a fully managed relational
engine, a cloud-native high-performance variant, high availability, read replicas, and a
migration service.

\begin{center}
\footnotesize
\begin{tabular}{@{}p{2.8cm}p{3.2cm}p{3.2cm}p{3.2cm}@{}}
\toprule
\textbf{Capability} & \textbf{AWS} & \textbf{Azure} & \textbf{GCP} \\
\midrule
Managed relational & Amazon RDS & Azure SQL Database & Cloud SQL \\
Cloud-native engine & Amazon Aurora & SQL Hyperscale & AlloyDB \\
Managed PostgreSQL & RDS / Aurora PostgreSQL & Azure DB for PostgreSQL & Cloud SQL / AlloyDB \\
High availability & Multi-AZ & Zone-redundant HA & Regional HA \\
Read scaling & Read Replicas & Read replicas / geo-replicas & Read Replicas \\
Global / geo & Aurora Global Database & Active geo-replication & Cross-region replicas \\
Serverless compute & Aurora Serverless v2 & Azure SQL Serverless & AlloyDB (elastic) \\
Migration service & DMS + SCT & Azure Database Migration Svc & Database Migration Svc \\
\bottomrule
\end{tabular}
\end{center}

\noindent \textbf{Snowflake}, \textbf{Fabric}, \textbf{Databricks}, and \textbf{MongoDB} are not
managed relational engines, but they solve overlapping problems: Snowflake and Fabric provide SQL
analytics over warehouse/lakehouse storage, Databricks consumes OLTP sources through Lakeflow
Connect, JDBC federation, and CDC landing into Delta tables, while MongoDB provides a distributed
document store whose schema design (Chapter~2) replaces rigid relational normalization with
intentional, access-driven modeling.

\section{Elastic Compute and Analytical Warehouses}
\index{elastic compute!equivalence}\index{virtual warehouse!equivalence}%
Snowflake's Chapter~2 topic---separating compute from storage and scaling it elastically---has a
counterpart in every analytics platform.

\begin{center}
\footnotesize
\begin{tabular}{@{}p{2.8cm}p{3.2cm}p{3.2cm}p{3.2cm}@{}}
\toprule
\textbf{Concept} & \textbf{Snowflake} & \textbf{AWS} & \textbf{Azure / GCP} \\
\midrule
Compute unit & Virtual Warehouse & Redshift RA3 / Serverless & Synapse pool / BigQuery slots \\
Elastic concurrency & Multi-cluster warehouse & Concurrency scaling & Autoscale / on-demand \\
Pause when idle & Auto-suspend / resume & Pause / resume & Auto-pause / serverless \\
Billing unit & Credits & RPU / node-hours & DWU / slot-hours \\
Result reuse & Result cache & Result cache & Result / BI cache \\
\bottomrule
\end{tabular}
\end{center}

\noindent \textbf{Microsoft Fabric} meters all workloads against a shared pool of Capacity Units
(CU) on an F-SKU, an approach closer to a single elastic budget than to per-engine warehouses.
\textbf{Databricks} meters compute in Databricks Units (DBUs) layered over the underlying cloud
VM cost: SQL warehouses (serverless, pro, or classic) play the virtual-warehouse role with
auto-stop and multi-cluster load balancing, while job clusters give per-run isolation.
\textbf{MongoDB Atlas} scales through cluster tiers with optional auto-scaling of compute and storage.

\section{Data Organization and the Medallion Pattern}
\index{medallion architecture!equivalence}%
Fabric's Chapter~2 topic---the Bronze/Silver/Gold medallion---is a portable design pattern rather
than a product. Every platform implements the same progressive-refinement idea with its own primitives.

\begin{center}
\footnotesize
\begin{tabular}{@{}p{2.6cm}p{3.4cm}p{3.4cm}p{3.0cm}@{}}
\toprule
\textbf{Layer} & \textbf{Fabric} & \textbf{Data lake (AWS/Azure/GCP)} & \textbf{Snowflake} \\
\midrule
Bronze (raw) & Bronze Lakehouse & \texttt{raw/} prefix (S3/ADLS/GCS) & RAW schema / stage \\
Silver (validated) & Silver Lakehouse & \texttt{validated/} prefix & STAGING schema \\
Gold (curated) & Gold Lakehouse & \texttt{curated/} prefix & ANALYTICS schema \\
Storage format & Delta-Parquet & Parquet / Delta / Iceberg & Micro-partitions \\
\bottomrule
\end{tabular}
\end{center}

\noindent \textbf{Databricks} is the medallion pattern's native home: \texttt{bronze},
\texttt{silver}, and \texttt{gold} schemas of Delta tables governed by Unity Catalog, with
expectations and quarantine as the quality gates between layers. \textbf{MongoDB} realizes the
same separation with distinct collections or databases (for example \texttt{raw\_events},
\texttt{validated}, \texttt{curated}) governed by schema-design patterns rather than a lake layout.

\section{Document Data Modeling}
\index{document modeling!equivalence}%
MongoDB's Chapter~2 schema-design patterns generalize to every document and wide-column store.

\begin{center}
\footnotesize
\begin{tabular}{@{}p{3.0cm}p{3.2cm}p{3.2cm}p{3.0cm}@{}}
\toprule
\textbf{Concept} & \textbf{MongoDB} & \textbf{Amazon DynamoDB} & \textbf{Azure Cosmos DB} \\
\midrule
Primary access unit & Document / collection & Item / table & Item / container \\
Partitioning & Shard key & Partition key & Partition key \\
One-table modeling & Embedding + patterns & Single-table design & Embedding + partitioning \\
Time-series grouping & Bucket Pattern & Composite sort keys & Bucketing by key \\
Optional attributes & Attribute Pattern & Sparse attributes & Sparse properties \\
\bottomrule
\end{tabular}
\end{center}

\noindent Google Cloud's closest analogs are \textbf{Firestore} and \textbf{Bigtable}. In
\textbf{Snowflake} and \textbf{Fabric}, semi-structured documents are stored in \texttt{VARIANT}
(Snowflake) or JSON columns and queried with path expressions rather than modeled as collections.
\textbf{Databricks} likewise has no document store; it ingests MongoDB and Cosmos DB through
Spark connectors and stores documents as Delta rows with \texttt{VARIANT}-style handling.

\section{Globally Distributed and NoSQL Databases}
\index{NoSQL!equivalence}\index{distributed databases!equivalence}%
Chapter~3 branches by platform into high-scale NoSQL (AWS DynamoDB, Azure Cosmos DB), globally
consistent NewSQL (GCP Cloud Spanner), and advanced document modeling (MongoDB). They all answer
the same question: how to scale a database horizontally across partitions and regions.

\begin{center}
\footnotesize
\begin{tabular}{@{}p{2.8cm}p{2.9cm}p{2.9cm}p{2.9cm}p{2.4cm}@{}}
\toprule
\textbf{Concept} & \textbf{DynamoDB} & \textbf{Cosmos DB} & \textbf{Cloud Spanner} & \textbf{MongoDB} \\
\midrule
Model & Key-value / doc & Multi-model & Distributed SQL & Document \\
Partitioning & Partition key & Partition key & Split by key range & Shard key \\
Throughput unit & RCU / WCU & Request Units (RU) & Nodes / processing units & Cluster tier \\
Secondary indexes & GSI / LSI & Automatic indexing & Secondary indexes & Indexes \\
Change stream & DynamoDB Streams & Change Feed & Change streams & Change streams \\
Expiry & TTL & TTL & (manual) & TTL index \\
Multi-region & Global Tables & Multi-region writes & Multi-region (synchronous) & Zones / global clusters \\
Consistency & Eventual / strong & 5 tunable levels & External (TrueTime) & Tunable read/write concern \\
\bottomrule
\end{tabular}
\end{center}

\noindent \textbf{Snowflake} and \textbf{Fabric} are analytical platforms rather than OLTP key-value
stores; they ingest from these operational databases and store semi-structured payloads in
\texttt{VARIANT}/JSON columns. Google's dedicated NoSQL analogs are \textbf{Firestore} (document)
and \textbf{Bigtable} (wide-column). \textbf{Databricks} is likewise analytical-only: it consumes
DynamoDB, Cosmos DB, and Bigtable through connectors into Delta rather than offering a
request-serving store of its own.

\section{Managed Apache Spark}
\index{Apache Spark!equivalence}%
Microsoft Fabric's Chapter~3 topic---managed Spark pools and notebooks---has a first-class
equivalent on every cloud.

\begin{center}
\footnotesize
\begin{tabular}{@{}p{2.8cm}p{3.2cm}p{3.2cm}p{3.2cm}@{}}
\toprule
\textbf{Capability} & \textbf{Fabric} & \textbf{AWS} & \textbf{Azure / GCP} \\
\midrule
Managed Spark & Starter / Custom pools & Glue / EMR / EMR Serverless & Synapse Spark / Databricks; Dataproc \\
Notebook UX & Fabric Notebooks & Glue Studio / EMR Studio & Synapse / Databricks; Vertex AI \\
Serverless option & Autoscaling pools & Glue / EMR Serverless & Synapse serverless; Dataproc Serverless \\
Lake format & Delta on OneLake & Delta / Iceberg / Hudi & Delta; Delta / Iceberg \\
\bottomrule
\end{tabular}
\end{center}

\noindent \textbf{Databricks} is the reference implementation this table generalizes: the
Databricks Runtime (with the Photon vectorized engine) over autoscaling or serverless clusters,
with Delta as the default table format. \textbf{Snowflake} offers \textbf{Snowpark} for
DataFrame-style processing, and \textbf{MongoDB} integrates through the Spark Connector.

\section{Query Acceleration and Materialized Views}
\index{materialized views!equivalence}\index{query acceleration!equivalence}%
Snowflake's Chapter~3 acceleration tools---materialized views, the Search Optimization Service, and
transient/temporary tables---map to comparable features in every warehouse.

\begin{center}
\footnotesize
\begin{tabular}{@{}p{2.8cm}p{3.0cm}p{3.0cm}p{3.4cm}@{}}
\toprule
\textbf{Feature} & \textbf{Snowflake} & \textbf{AWS Redshift} & \textbf{Azure Synapse / GCP BigQuery} \\
\midrule
Materialized views & Materialized Views & Materialized views & Materialized views \\
Point-lookup accel. & Search Optimization Svc & Sort keys / zone maps & (indexing) / search indexes \\
Result reuse & Result cache & Result cache & Result-set cache / BI Engine \\
Ephemeral tables & Transient / temporary & Temp tables & Temp tables \\
\bottomrule
\end{tabular}
\end{center}

\noindent \textbf{Fabric} accelerates reads through Direct Lake and its Warehouse engine,
\textbf{Databricks} through materialized views, Photon, file-level statistics with bloom filters,
and predictive optimization over Delta tables, while \textbf{MongoDB} relies on secondary and
compound indexes plus the in-memory working set.

\section{In-Memory Caching and Hybrid Ingestion}
\index{caching!equivalence}\index{in-memory cache!equivalence}%
The AWS and Azure Chapter~4 topic---low-latency caching in front of a durable store, fed by hybrid
(edge-to-cloud) ingestion---has a managed Redis-compatible equivalent on every cloud.

\begin{center}
\footnotesize
\begin{tabular}{@{}p{2.8cm}p{3.0cm}p{3.0cm}p{3.4cm}@{}}
\toprule
\textbf{Capability} & \textbf{AWS} & \textbf{Azure} & \textbf{GCP / others} \\
\midrule
Managed cache & ElastiCache (Redis/Memcached) & Cache for Redis & Memorystore (Redis/Memcached) \\
Real-time API layer & AppSync / API Gateway & API Management / SignalR & API Gateway / Apigee \\
Hybrid ingestion & DataSync / Snow family / IoT Core & Data Box / IoT Hub / Arc & Transfer Appliance / IoT Core \\
Streaming intake & Kinesis / MSK & Event Hubs & Pub/Sub \\
\bottomrule
\end{tabular}
\end{center}

\noindent \textbf{Snowflake} and \textbf{Fabric} rely on result and Direct~Lake caches rather than a
standalone Redis tier, \textbf{Databricks} relies on its warehouse result cache and the
cluster-local disk cache rather than a general-purpose cache service, while \textbf{MongoDB}
serves hot data from its in-memory working set and Atlas can pair with an external cache.

\section{Wide-Column and Key-Value NoSQL}
\index{wide-column!equivalence}\index{Bigtable!equivalence}%
GCP's Chapter~4 topic---Cloud Bigtable---is a wide-column, high-throughput, low-latency NoSQL store.
Its peers span the managed and open-source ecosystems.

\begin{center}
\footnotesize
\begin{tabular}{@{}p{2.8cm}p{3.0cm}p{3.0cm}p{3.4cm}@{}}
\toprule
\textbf{Characteristic} & \textbf{GCP Bigtable} & \textbf{AWS} & \textbf{Azure / open source} \\
\midrule
Data model & Wide-column & DynamoDB (key-value/doc) & Cosmos DB (multi-model) \\
HBase-compatible & HBase API & Keyspaces (Cassandra) & Managed Instance for Cassandra \\
Scaling unit & Nodes / clusters & On-demand / provisioned & RU/s (Cosmos) \\
Best fit & Time-series, IoT, ad-tech & Serverless key-value & Global multi-model \\
\bottomrule
\end{tabular}
\end{center}

\noindent \textbf{Snowflake} and \textbf{Fabric} address these workloads analytically rather than as
operational key-value stores; \textbf{Databricks} reads Cassandra- and HBase-family sources with
Spark connectors for analytical offload; \textbf{MongoDB} covers the document end of the same
NoSQL spectrum.

\section{Data Protection: Cloning, Time Travel, and Migration}
\index{cloning!equivalence}\index{time travel!equivalence}\index{point-in-time recovery!equivalence}%
Snowflake's Chapter~4 topic---zero-copy cloning, Time Travel, and Fail-safe---plus MongoDB's
relational-to-document migration map to data-protection and migration tooling across the clouds.

\begin{center}
\footnotesize
\begin{tabular}{@{}p{2.8cm}p{3.0cm}p{3.0cm}p{3.4cm}@{}}
\toprule
\textbf{Capability} & \textbf{Snowflake} & \textbf{AWS / Azure} & \textbf{GCP / Fabric} \\
\midrule
Instant clone & Zero-copy clone & Aurora clone / SQL DB copy & BigQuery table clone; OneLake shortcut \\
Point-in-time recovery & Time Travel & PITR (RDS/SQL DB backups) & PITR; Delta time travel \\
Retention safety net & Fail-safe (7 days) & Automated backups & Backups; Delta VACUUM window \\
Schema migration & --- & DMS / SCT & Database Migration Service \\
\bottomrule
\end{tabular}
\end{center}

\noindent \textbf{MongoDB} migrations use \textbf{Relational Migrator} plus the embed-vs-reference
patterns from Chapter~4, \textbf{Delta Lake} exposes historical versions through
\texttt{VERSION AS OF} time-travel queries, and \textbf{Databricks} rounds out the set with
\texttt{RESTORE}, \texttt{VACUUM}, and zero-copy \emph{shallow clones} of Delta tables---the
closest analog to Snowflake's zero-copy cloning outside the warehouse.

\section{Account and Environment Setup}
\index{account setup!equivalence}\index{cost guardrails!equivalence}%
Chapter~0's first-hour pattern---free offer, secured administrative identity, CLI, cost
guardrails, verification---repeats on every platform, with different defaults for what a
``free'' environment means.

\begin{center}
\footnotesize
\begin{tabular}{@{}p{2.4cm}p{2.9cm}p{2.9cm}p{2.9cm}p{2.9cm}@{}}
\toprule
\textbf{Setup step} & \textbf{AWS} & \textbf{Azure} & \textbf{GCP} & \textbf{Databricks} \\
\midrule
Free offer & Free Tier (12 mo + always free) & Free account (\$200/30 days) & Free trial (\$300/90 days) & 14-day trial / Free Edition \\
Admin identity & Root (MFA, then never) & Global Admin (avoid) & Organization Owner (avoid) & Account admin (first user) \\
Workforce identity & IAM Identity Center & Microsoft Entra ID & Cloud Identity & SSO + SCIM from IdP \\
CLI & AWS CLI v2 (\texttt{aws configure}) & Azure CLI (\texttt{az login}) & gcloud (\texttt{gcloud init}) & \texttt{databricks configure} \\
Cost guardrail & AWS Budgets + alerts & Cost Management budgets & Budgets \& alerts & Budget policies + serverless guardrails \\
Browser shell & CloudShell & Cloud Shell & Cloud Shell & Web terminal / notebooks \\
\bottomrule
\end{tabular}
\end{center}

\noindent \textbf{Fabric} starts from a Microsoft 365/Power BI work account plus a trial capacity
(no card, org tenant); \textbf{Snowflake} uses a 30-day/\$400 trial guarded by Resource Monitors;
\textbf{MongoDB Atlas} offers a forever-free M0 cluster secured by database users and an IP
allowlist.

\section{Cloud Data Warehouses and MPP Architecture}
\index{data warehouse!equivalence}\index{MPP!equivalence}%
Chapter~5 introduces the cloud data warehouse. The MPP mental model---shared-nothing compute over
distributed columnar storage---transfers to every platform, even where the service hides it.

\begin{center}
\footnotesize
\begin{tabular}{@{}p{2.4cm}p{2.9cm}p{2.9cm}p{2.9cm}p{2.9cm}@{}}
\toprule
\textbf{Capability} & \textbf{AWS} & \textbf{Azure} & \textbf{GCP} & \textbf{Databricks} \\
\midrule
Cloud data warehouse & Amazon Redshift & Synapse dedicated pools / Fabric Warehouse & BigQuery & Databricks SQL warehouses \\
Serverless option & Redshift Serverless & Fabric capacity / Synapse serverless & BigQuery (default) & Serverless SQL warehouses \\
Managed storage layer & RA3 managed storage & ADLS-backed separation & Capacitor on Colossus & Delta on your object storage \\
Bulk load & COPY & COPY INTO / bulk insert & LOAD DATA / \texttt{bq load} & \texttt{COPY INTO} / Auto Loader \\
Change-data upsert & MERGE & MERGE & MERGE & \texttt{MERGE INTO} \\
BI acceleration & Materialized views & Materialized views / semantic models & MVs, BI Engine & Materialized views, Photon \\
\addlinespace
Pricing lever (provisioned) & Node-hours + storage TB & DWU-hours + storage & Flat-rate slots & DBU + VM cost (classic) \\
Pricing lever (serverless) & RPU-hours & Capacity units & Bytes scanned / slots & Serverless DBU-hours \\
Concurrency burst & Concurrency scaling & Workload groups & Autoscaling slots & Multi-cluster load balancing \\
\bottomrule
\end{tabular}
\end{center}

\noindent \textbf{Fabric} blurs the warehouse/lake boundary: warehouse and lakehouse share OneLake
Delta--Parquet storage. \textbf{Snowflake} separates storage, virtual warehouses, and cloud
services into three independently scaled layers---the closest conceptual cousin to RA3.
\textbf{MongoDB} is an operational document store, not a warehouse; the contrast explains why
star schemas live outside the system of record.

\section{Warehouse Performance and Physical Design}
\index{performance tuning!equivalence}\index{distribution style!equivalence}%
Chapter~6's physical-design levers---move less data, read less data, schedule work fairly---exist
on every analytical engine; each platform exposes a different subset.

\begin{center}
\footnotesize
\begin{tabular}{@{}p{2.4cm}p{2.9cm}p{2.9cm}p{2.9cm}p{2.9cm}@{}}
\toprule
\textbf{Capability} & \textbf{AWS} & \textbf{Azure} & \textbf{GCP} & \textbf{Databricks} \\
\midrule
Data distribution & DISTSTYLE EVEN/KEY/ALL & Synapse hash/round-robin/replicated & Automatic & Automatic (partitioning optional) \\
Sort optimization & SORTKEY (compound/interleaved) & Clustered columnstore index & Clustering columns & Z-ORDER / liquid clustering \\
Pre-computation & Materialized views & Materialized views & Materialized views & Materialized views \\
Workload governance & WLM + QMRs & Workload groups / classifiers & Reservations \& slots & SQL warehouse clusters + queries API \\
Burst concurrency & Concurrency scaling & Scale-out / serverless burst & Autoscaling slots & Multi-cluster autoscaling \\
\addlinespace
Manual tuning burden & Medium (AUTO helps) & Medium-high & Low & Low (predictive optimization) \\
Skew diagnosis & SVV\_TABLE\_INFO & DMVs & Query plan explanation & Query profile, AQE skew split \\
\bottomrule
\end{tabular}
\end{center}

\noindent \textbf{Snowflake} removes most physical levers---micro-partition pruning is automatic
and clustering keys optional---trading control for operational simplicity. \textbf{Fabric}
inherits SQL Server-style distribution choices on dedicated capacity. \textbf{MongoDB} offers
indexes and shard keys rather than distribution styles, because its workload is point lookups,
not scans.

\section{Lakehouse Query Layer and Open Table Formats}
\index{lakehouse!equivalence}\index{table format!equivalence}%
Chapter~7's lakehouse pattern---open files on object storage, a shared catalog, multiple query
engines---exists on every cloud, with Databricks as its original author.

\begin{center}
\footnotesize
\begin{tabular}{@{}p{2.4cm}p{2.9cm}p{2.9cm}p{2.9cm}p{2.9cm}@{}}
\toprule
\textbf{Capability} & \textbf{AWS} & \textbf{Azure} & \textbf{GCP} & \textbf{Databricks} \\
\midrule
Serverless SQL on the lake & Athena & Synapse/Fabric serverless SQL & BigQuery external tables & SQL warehouses over Delta \\
Shared metastore & Glue Data Catalog & Purview / OneLake catalog & Dataplex / BigLake & Unity Catalog \\
Schema inference & Glue crawlers & Synapse auto-discovery & BigQuery autodetect & Auto Loader inference + hints \\
Warehouse-to-lake query & Redshift Spectrum & Lakehouse shortcuts & BigLake tables & Native (same tables) \\
Table format & Iceberg / Hudi / Delta on S3 & Delta (OneLake native) & BigLake Iceberg & Delta (UniForm for Iceberg/Hudi reads) \\
\addlinespace
Pricing lever & Per TB scanned & Capacity units & Bytes scanned / slots & DBU-hours \\
Cost control & Partitioning + Parquet & Partitioning + Delta & Partitioning + clustering & Liquid clustering + OPTIMIZE \\
\bottomrule
\end{tabular}
\end{center}

\noindent \textbf{Fabric} is the most opinionated implementation: OneLake plus shortcuts makes the
warehouse and lake the same physical data. \textbf{Snowflake} reads external Iceberg tables while
its native tables remain the managed counterpart. \textbf{MongoDB} plays the source-system role;
Atlas Data Federation queries S3 but the document model stays operational.

\section{Lake Governance, Sharing, and Clean Rooms}
\index{lake governance!equivalence}\index{clean room!equivalence}%
Chapter~8's governance primitives---catalog, tag-based policy, fine-grained controls, clean-room
collaboration---have converged across platforms.

\begin{center}
\footnotesize
\begin{tabular}{@{}p{2.4cm}p{2.9cm}p{2.9cm}p{2.9cm}p{2.9cm}@{}}
\toprule
\textbf{Capability} & \textbf{AWS} & \textbf{Azure} & \textbf{GCP} & \textbf{Databricks} \\
\midrule
Lake governance layer & Lake Formation & Microsoft Purview & Dataplex & Unity Catalog \\
Tag-based access & LF-Tags & Purview classifications & Policy tags & UC tags + attribute-based policies \\
Row/column security & LF data filters + cell filters & SQL RLS/masking & BigQuery RLS + column ACLs & Row filters + column masks \\
Live cross-account share & Redshift data sharing & Fabric OneLake sharing & Analytics Hub & Delta Sharing (open protocol) \\
Managed clean room & AWS Clean Rooms & Azure confidential clean rooms & BigQuery data clean rooms & Databricks Clean Rooms \\
\addlinespace
Enforcement point & Catalog (engine-agnostic) & Per engine + Purview & BigQuery policy tags & UC (all engines on platform) \\
Default stance & IAMAllowedPrincipals (revoke!) & Deny until granted & Deny until granted & Deny until granted \\
\bottomrule
\end{tabular}
\end{center}

\noindent \textbf{Snowflake} implements the same primitives natively---row access policies,
masking policies, tags, zero-copy shares/listings---with a clean-room offering parallel to AWS
Clean Rooms. \textbf{MongoDB} enforces field-level security at the document layer (client-side
field-level encryption). Delta Sharing's distinction is the \emph{open} protocol: recipients read
with pandas, Spark, or Power BI without any Databricks account.

\section{Managed Hadoop and Spark Clusters}
\index{EMR!equivalence}\index{managed Spark!equivalence}%
Chapter~9's deployment-mode spectrum---full-control clusters to fully serverless Spark---is the
portable idea; every cloud occupies several points on it.

\begin{center}
\footnotesize
\begin{tabular}{@{}p{2.4cm}p{2.9cm}p{2.9cm}p{2.9cm}p{2.9cm}@{}}
\toprule
\textbf{Capability} & \textbf{AWS} & \textbf{Azure} & \textbf{GCP} & \textbf{Databricks} \\
\midrule
Managed Hadoop/Spark & EMR on EC2 & HDInsight & Dataproc & All-purpose clusters (no YARN) \\
Serverless Spark & EMR Serverless & Fabric / Synapse Spark & Dataproc Serverless & Serverless compute \\
Spark on Kubernetes & EMR on EKS & AKS-hosted Spark & Dataproc on GKE & (control plane managed) \\
Object-storage filesystem & EMRFS (S3) & ABFS (ADLS) & GCS connector & DBFS mounts / UC volumes \\
Cheap interruptible nodes & Spot task nodes & Spot node pools & Preemptible workers & Spot with fallback policy \\
\addlinespace
Cluster lifetime model & Transient vs.\ long-running & Both & Both & All-purpose vs.\ job clusters \\
Notebook integration & EMR Studio & Fabric notebooks & Vertex AI / Jupyter & Databricks notebooks (native) \\
\bottomrule
\end{tabular}
\end{center}

\noindent \textbf{Fabric} folds Spark into its SaaS lakehouse---pools and notebooks, never node
types. \textbf{Databricks} popularized the transient, job-scoped cluster pattern this chapter
teaches; cluster policies replace EMR's instance-fleet discipline. \textbf{Snowflake}'s Snowpark
pushes DataFrame code into the warehouse instead of a cluster. \textbf{MongoDB} stays the
operational source Spark jobs read via the connector.

\section{Serverless, Catalog-Integrated ETL}
\index{ETL!equivalence}\index{serverless ETL!equivalence}%
Chapter~10's serverless ETL pattern---schema-aware transforms over a shared catalog, billed per
run---is universal; only the branding differs.

\begin{center}
\footnotesize
\begin{tabular}{@{}p{2.4cm}p{2.9cm}p{2.9cm}p{2.9cm}p{2.9cm}@{}}
\toprule
\textbf{Capability} & \textbf{AWS} & \textbf{Azure} & \textbf{GCP} & \textbf{Databricks} \\
\midrule
Serverless ETL engine & AWS Glue (Spark) & Fabric Dataflow Gen2 / pipelines & Dataflow / Data Fusion & Lakeflow Declarative Pipelines (DLT) \\
Visual designer & Glue Studio & Data Factory designer & Data Fusion studio & DLT graph editor / notebooks \\
Interactive data prep & Glue DataBrew & Power Query (Fabric) & Dataprep & Databricks notebooks + Assistant \\
Shared catalog & Glue Data Catalog & OneLake catalog / Purview & Dataplex & Unity Catalog \\
Re-run state & Job bookmarks & Pipeline state / watermarks & Job state / watermarking & Checkpoints + pipeline state \\
\addlinespace
Capacity unit & DPU-hour & Capacity units & vCPU-hr & DBU-hour (serverless or classic) \\
Spot discount & Flex execution & Low-priority pools & Preemptible/Spot workers & Spot instances with fallback \\
\bottomrule
\end{tabular}
\end{center}

\noindent \textbf{Snowflake} approaches the same problems from inside the warehouse: streams,
tasks, and Snowpark replace external ETL for in-warehouse transforms. \textbf{MongoDB} users meet
this layer as the Atlas Data API or Spark connector feeding the lake. DLT's distinguishing move
is \emph{declarative} pipelines with expectations (warn/drop/fail) as first-class quality gates.

\section{Change Data Capture, Event Routing, and Schema Contracts}
\index{CDC!equivalence}\index{schema registry!equivalence}%
Chapter~11's disciplines---CDC replication, event routing, and schema contracts---are
cloud-agnostic; each platform offers a managed shape of the same ideas.

\begin{center}
\footnotesize
\begin{tabular}{@{}p{2.4cm}p{2.9cm}p{2.9cm}p{2.9cm}p{2.9cm}@{}}
\toprule
\textbf{Capability} & \textbf{AWS} & \textbf{Azure} & \textbf{GCP} & \textbf{Databricks} \\
\midrule
Managed CDC replication & AWS DMS & Azure DMS / Data Factory CDC & Datastream & Lakeflow Connect (managed connectors) \\
Event router & EventBridge & Event Grid & Eventarc & (cloud-native router triggers jobs) \\
Serverless functions & Lambda & Azure Functions & Cloud Run functions & (SQL warehouses + jobs; cloud functions alongside) \\
Schema registry & Glue Schema Registry & Azure Schema Registry & Managed Service for Schema Registry & Schema inference + UC schema enforcement \\
Contract enforcement & In Glue jobs / registry checks & In pipelines & In Dataflow / registry checks & Delta constraints + DLT expectations \\
\addlinespace
CDC lag metric & CDCLatencySource/Target & Replication latency & Datastream freshness & Pipeline event log freshness \\
CDC apply & Glue/EMR MERGE & T-SQL MERGE & BigQuery MERGE & \texttt{AUTO CDC} / \texttt{MERGE INTO} \\
\bottomrule
\end{tabular}
\end{center}

\noindent \textbf{Snowflake} ingests CDC through Snowpipe and streams, applying contracts at the
table boundary. \textbf{MongoDB} contributes change streams---the document-native CDC primitive
that pairs with everything here. \textbf{Fabric} routes database mirroring and event streams
through OneLake, its version of the CDC-to-lake path.

\section{Orchestration: DAG Schedulers and State Machines}
\index{orchestration!equivalence}\index{Airflow!equivalence}%
Chapter~12 splits orchestration into two families---general-purpose DAG schedulers and native
state machines; both exist on every cloud.

\begin{center}
\footnotesize
\begin{tabular}{@{}p{2.4cm}p{2.9cm}p{2.9cm}p{2.9cm}p{2.9cm}@{}}
\toprule
\textbf{Capability} & \textbf{AWS} & \textbf{Azure} & \textbf{GCP} & \textbf{Databricks} \\
\midrule
Managed Airflow & Amazon MWAA & Data Factory managed Airflow & Cloud Composer & (external Airflow via Jobs API) \\
Native state machine & Step Functions & Logic Apps / Durable Functions & Workflows & Workflows (jobs + tasks) \\
File-arrival trigger & EventBridge + sensor & Event Grid trigger & Eventarc trigger & File-arrival triggers \\
Failure alerting & SNS via operator/rule & Monitor alerts & Alerting to Pub/Sub & Job notifications + system tables \\
\addlinespace
Inter-task data & XComs (small) / S3 (large) & Pipeline variables / ADLS & XComs / GCS & Task values / Delta tables \\
Partial recovery & Redrive / Airflow retries & Rerun from activity & Step retry blocks & Repair run (failed + downstream) \\
\bottomrule
\end{tabular}
\end{center}

\noindent \textbf{Fabric} pipelines play the orchestrator role inside the SaaS boundary.
\textbf{Snowflake} tasks and DAGs orchestrate in-warehouse work, and dbt orchestrates
transformation graphs at the SQL layer. \textbf{MongoDB} triggers and Atlas schedulers cover the
document-store corner. Workflows' native advantage is colocation: clusters, catalog, and lineage
in one pane.

\section{Messaging and Streaming Primitives}
\index{streaming!equivalence}\index{messaging!equivalence}%
Chapter~13's three primitives---queues, pub/sub, append-only logs---ship on every cloud.

\begin{center}
\footnotesize
\begin{tabular}{@{}p{2.4cm}p{2.9cm}p{2.9cm}p{2.9cm}p{2.9cm}@{}}
\toprule
\textbf{Capability} & \textbf{AWS} & \textbf{Azure} & \textbf{GCP} & \textbf{Databricks} \\
\midrule
Pull queue & Amazon SQS & Queue Storage / Service Bus & Pub/Sub (pull) & (none; consumes external queues) \\
Pub/sub fan-out & Amazon SNS & Event Grid / Service Bus topics & Pub/Sub & (none; consumes external buses) \\
Durable streaming log & Kinesis Data Streams & Event Hubs & Pub/Sub (ordered) & Kafka sources into Structured Streaming \\
Managed Kafka & Amazon MSK & Event Hubs for Kafka & Managed Service for Apache Kafka & (Consumes MSK/Confluent/Event Hubs) \\
Serverless stream processor & Managed Flink / Lambda & Stream Analytics & Dataflow & Structured Streaming (serverless option) \\
\addlinespace
Scaling unit & Shard (1 MB/s in, 2 MB/s out) & Throughput units / partitions & Throughput quotas & Kafka partitions / trigger rate \\
Ordering scope & Per partition key & Per partition key & Per ordering key & Per Kafka partition \\
Retention & 24 h--365 days & 1--90 days & Up to 31 days & Source-side (Kafka/Kinesis retention) \\
\bottomrule
\end{tabular}
\end{center}

\noindent \textbf{Fabric} Real-Time Intelligence wraps KQL databases and Eventstreams around the
same log concept. \textbf{Snowflake} consumes streams via Snowpipe Streaming. \textbf{MongoDB}
change streams are the database-native analog---replayable, resumable. Databricks deliberately
ships no bus: it is the consumer/processor of record for all of them.

\section{Managed Stream-to-Sink Delivery}
\index{stream delivery!equivalence}\index{Firehose!equivalence}%
Chapter~14's buffer--transform--convert--land delivery pattern exists everywhere, usually as a
feature of the streaming service rather than a separate product.

\begin{center}
\footnotesize
\begin{tabular}{@{}p{2.4cm}p{2.9cm}p{2.9cm}p{2.9cm}p{2.9cm}@{}}
\toprule
\textbf{Capability} & \textbf{AWS} & \textbf{Azure} & \textbf{GCP} & \textbf{Databricks} \\
\midrule
Managed stream delivery & Data Firehose & Event Hubs Capture & Pub/Sub export / Dataflow templates & Auto Loader + Structured Streaming \\
Inline transformation & Lambda in Firehose & Functions pattern & Dataflow / Cloud Run functions & Streaming DataFrame transforms \\
In-flight format conversion & JSON $\rightarrow$ Parquet/ORC & Via downstream Synapse/Fabric & Via Dataflow & Native Delta writes \\
Warehouse delivery & Firehose $\rightarrow$ Redshift & Event Hubs $\rightarrow$ Synapse & Pub/Sub $\rightarrow$ BigQuery & Streaming $\rightarrow$ Delta tables \\
Search delivery & Firehose $\rightarrow$ OpenSearch & Event Hubs $\rightarrow$ Elastic & Pub/Sub $\rightarrow$ Elasticsearch & Sink connectors / foreachBatch \\
\addlinespace
Buffering model & Size + interval hints & Time/size capture windows & Subscription batching & Micro-batch triggers \\
Failure handling & S3 backup + error prefixes & Dead-lettering & Dead-letter topics & \texttt{\_rescued\_data} + checkpoint replay \\
\bottomrule
\end{tabular}
\end{center}

\noindent \textbf{Snowflake} Snowpipe Streaming is the same idea aimed at one destination: the
warehouse table. \textbf{Fabric} Eventstreams route to lakehouse tables with optional
transformation. \textbf{MongoDB} sits on the receiving end via sink connectors or trigger-based
loads. On Databricks the ``delivery service'' is just a streaming job---more control, more to
operate.

\section{Stateful Stream Processing}
\index{stream processing!equivalence}\index{Flink!equivalence}%
Chapter~15 is the deepest discipline in the book: event time, watermarks, windows, state, and
checkpoints are identical everywhere; only the managed packaging differs.

\begin{center}
\footnotesize
\begin{tabular}{@{}p{2.4cm}p{2.9cm}p{2.9cm}p{2.9cm}p{2.9cm}@{}}
\toprule
\textbf{Capability} & \textbf{AWS} & \textbf{Azure} & \textbf{GCP} & \textbf{Databricks} \\
\midrule
Engine & Managed Service for Apache Flink & HDInsight/AKS-hosted Flink & Dataflow (Beam model) & Structured Streaming (Spark) \\
Execution model & True streaming & True streaming & Unified batch+stream & Micro-batch (continuous optional) \\
Native stream SQL & Flink SQL & Stream Analytics query language & Dataflow SQL / Beam & Spark SQL over streams \\
Event-time windows & TUMBLE/HOP/SESSION & ASA windowing & Beam windows & \texttt{window()}, session windows \\
Lateness control & Watermarks + allowed lateness & ASA late-arrival policy & Beam watermarks & \texttt{withWatermark} \\
Custom state & Keyed process functions & Per platform & Beam stateful DoFns & \texttt{flatMapGroupsWithState} \\
\addlinespace
State backend & RocksDB (managed) & Per platform & Per runner & HDFS-compatible / RocksDB store \\
Exactly-once sink & Two-phase commit connectors & Idempotent sinks & Exactly-once in Dataflow & Checkpoint + idempotent Delta sink \\
\bottomrule
\end{tabular}
\end{center}

\noindent \textbf{Fabric} Real-Time Intelligence covers many windowed cases with KQL, trading
control for simplicity. \textbf{Snowflake} approaches streaming with Snowpipe Streaming plus
dynamic tables. \textbf{MongoDB} stream processing (Atlas) applies the same windowed ideas to
change streams. Learn Flink's model once; everywhere else is a dialect---Spark's dialect is the
micro-batch trigger and the checkpoint directory.

\section{Real-Time Serving, BI, and Reverse ETL}
\index{BI!equivalence}\index{reverse ETL!equivalence}%
Chapter~16's serving layer---dashboards for humans, operational pushes back to business
systems---closes the analytics loop on every platform.

\begin{center}
\footnotesize
\begin{tabular}{@{}p{2.4cm}p{2.9cm}p{2.9cm}p{2.9cm}p{2.9cm}@{}}
\toprule
\textbf{Capability} & \textbf{AWS} & \textbf{Azure} & \textbf{GCP} & \textbf{Databricks} \\
\midrule
Managed BI & Amazon QuickSight & Power BI (Fabric) & Looker / Looker Studio & AI/BI dashboards \\
In-memory acceleration & SPICE & Import / semantic models & BI Engine & SQL warehouse result cache \\
NL analytics & Amazon Q in QuickSight & Copilot (Power BI) & Gemini in Looker & Genie spaces \\
Row-level security & QuickSight RLS rules & Power BI RLS & Looker access filters & UC row filters (inherited) \\
Warehouse export & Redshift UNLOAD & COPY INTO / export & BigQuery EXPORT DATA & \texttt{COPY INTO} / Delta Sharing \\
Reverse-ETL glue & Lambda + API calls & Logic Apps / Functions & Cloud Run functions & Jobs + webhooks / partner tools \\
\addlinespace
Reader pricing & Per-reader session & Per-user / capacity & Per-user & Included in platform \\
\bottomrule
\end{tabular}
\end{center}

\noindent \textbf{Snowflake} hosts dashboards in Snowsight and streams reverse ETL through
external functions; the ecosystem (Census, Hightouch) serves all platforms. \textbf{MongoDB}
Charts covers operational BI directly on collections. AI/BI's differentiator: tiles read governed
Delta tables in place and inherit Unity Catalog lineage back to the producing pipeline.

\section{In-Warehouse Machine Learning}
\index{in-warehouse ML!equivalence}\index{Redshift ML!equivalence}%
Chapter~17's pattern---bringing the model to the data instead of shipping data to a modeling
platform---is now offered by every warehouse.

\begin{center}
\footnotesize
\begin{tabular}{@{}p{2.4cm}p{2.9cm}p{2.9cm}p{2.9cm}p{2.9cm}@{}}
\toprule
\textbf{Capability} & \textbf{AWS} & \textbf{Azure} & \textbf{GCP} & \textbf{Databricks} \\
\midrule
SQL-based ML & Redshift ML (CREATE MODEL) & Fabric SynapseML / T-SQL PREDICT & BigQuery ML (CREATE MODEL) & SQL functions over served models (\texttt{ai\_query}) \\
AutoML behind the scenes & SageMaker Autopilot & Azure ML AutoML & Vertex AI AutoML & AutoML (glass-box notebooks) \\
In-warehouse scoring & SQL function per model & T-SQL PREDICT & ML.PREDICT & Batch scoring UDFs from registry \\
Custom code path & BYO model via SageMaker & Synapse Spark MLlib & Remote models / Vertex & Notebooks + MLflow (native) \\
Clustering & K-Means problem type & MLlib KMeans & ML.KMEANS & MLlib / scikit-learn on Spark \\
\addlinespace
Training data leaves warehouse? & Yes (to SageMaker, in-account) & Varies & No (BigQuery-native) & No (Delta tables stay governed) \\
Cost driver & SageMaker training + S3 staging & Capacity units & Bytes processed + training & Cluster DBUs \\
\bottomrule
\end{tabular}
\end{center}

\noindent \textbf{Snowflake} ML offers \texttt{CREATE MODEL} and Snowpark ML with warehouse-side
scoring---the closest analog. \textbf{MongoDB} has no in-database ML; Atlas Vector Search
addresses the adjacent similarity problem. Databricks inverts the direction: the lakehouse
\emph{is} the modeling platform, so ``in-warehouse ML'' is the default rather than a feature.

\section{Managed ML Platforms}
\index{ML platform!equivalence}\index{SageMaker!equivalence}%
Chapter~18's managed-ML shape---a studio, managed training, a registry, deployment targets---maps
name-for-name across platforms.

\begin{center}
\footnotesize
\begin{tabular}{@{}p{2.4cm}p{2.9cm}p{2.9cm}p{2.9cm}p{2.9cm}@{}}
\toprule
\textbf{Capability} & \textbf{AWS} & \textbf{Azure} & \textbf{GCP} & \textbf{Databricks} \\
\midrule
ML studio & SageMaker Studio & Azure ML studio & Vertex AI Workbench & Databricks workspace + notebooks \\
Interactive data prep & Data Wrangler & Fabric Data Wrangler / Power Query & BigQuery DataFrames / Colab & Notebooks + Assistant \\
AutoML & SageMaker Autopilot & Azure ML AutoML & Vertex AI AutoML & Databricks AutoML \\
Managed training & SageMaker Training Jobs & Azure ML jobs & Vertex AI Training & Jobs on ML runtimes \\
Real-time inference & Real-time endpoints & Managed online endpoints & Vertex AI endpoints & Model Serving (serverless) \\
Batch inference & Batch Transform & Batch endpoints & Batch prediction & Spark UDF from registry + jobs \\
\addlinespace
Custom code & BYO container (any framework) & Custom environments & Custom containers & Custom images / libraries \\
GPU access & ml.g/ml.p instances & GPU compute clusters & GPU nodes/Vertex & GPU clusters (ML runtime) \\
\bottomrule
\end{tabular}
\end{center}

\noindent \textbf{Fabric} data science runs in-notebook with SynapseML against OneLake.
\textbf{Snowflake} ML covers training and serving inside the warehouse boundary.
\textbf{MongoDB} feeds these platforms as a source; its vector search serves retrieval rather
than training.

\section{MLOps: Feature Stores, Registries, and Monitoring}
\index{MLOps!equivalence}\index{feature store!equivalence}%
Chapter~19's MLOps primitives---feature stores, pipeline orchestration, registries, drift
monitoring---are platform-standard.

\begin{center}
\footnotesize
\begin{tabular}{@{}p{2.4cm}p{2.9cm}p{2.9cm}p{2.9cm}p{2.9cm}@{}}
\toprule
\textbf{Capability} & \textbf{AWS} & \textbf{Azure} & \textbf{GCP} & \textbf{Databricks} \\
\midrule
Feature store & SageMaker Feature Store & Fabric feature store / Feathr & Vertex AI Feature Store & Feature Engineering in UC \\
ML pipelines & SageMaker Pipelines & Azure ML pipelines & Vertex AI Pipelines (KFP) & Workflows + DLT / MLflow Projects \\
Model registry & SageMaker Model Registry & Azure ML registry & Vertex AI Model Registry & Unity Catalog model registry \\
Drift monitoring & SageMaker Model Monitor & Azure ML monitoring & Vertex AI Model Monitoring & Lakehouse Monitoring \\
Retrain trigger & EventBridge + pipeline & Logic Apps + pipeline & Eventarc + pipeline & Job triggers + webhooks \\
\addlinespace
Online store backing & DynamoDB-backed records & Cosmos-backed & Bigtable-backed & Feature serving endpoints \\
Lineage & SageMaker Lineage + OpenLineage & MLflow tracking & Vertex ML Metadata & Automatic UC lineage (data $\rightarrow$ model) \\
\bottomrule
\end{tabular}
\end{center}

\noindent \textbf{Snowflake}'s Feature Store and Model Registry bring the same lifecycle inside
the warehouse. \textbf{MongoDB} typically appears as the online-serving store behind low-latency
feature lookups in self-managed stacks. \textbf{MLflow}---born at Databricks---is the open-source
reference most of these registries interoperate with.

\section{Pre-trained AI Services and Unstructured Data}
\index{AI services!equivalence}\index{OCR!equivalence}%
Chapter~20's commodity building blocks---vision, OCR, NLP, translation---differ only in packaging;
the skill is pipeline integration, not model building.

\begin{center}
\footnotesize
\begin{tabular}{@{}p{2.4cm}p{2.9cm}p{2.9cm}p{2.9cm}p{2.9cm}@{}}
\toprule
\textbf{Capability} & \textbf{AWS} & \textbf{Azure} & \textbf{GCP} & \textbf{Databricks} \\
\midrule
OCR / document parsing & Textract & Document Intelligence & Document AI & \texttt{ai\_parse\_document} / partner parsers \\
Image/video labeling & Rekognition & Azure AI Vision & Vision AI / Video AI & Foundation Model APIs (vision models) \\
NLP (entities, sentiment, PII) & Comprehend & Azure AI Language & Natural Language AI & \texttt{ai\_query} with LLM endpoints \\
Translation & Translate & Translator & Translation API & \texttt{ai\_query} / external models \\
Foundation models & Amazon Bedrock & Azure OpenAI & Vertex AI & Foundation Model APIs (pay-per-token) \\
Vector index / search & OpenSearch / Kendra & AI Search & Vertex AI Vector Search & Mosaic AI Vector Search \\
\addlinespace
PII detection & Comprehend PII + Macie & Language PII + Purview & Sensitive Data Protection & UC discovery + \texttt{ai\_query} redaction \\
RAG orchestration & Bedrock Knowledge Bases & Azure AI RAG pattern & Vertex RAG Engine & LangChain + Vector Search + MLflow \\
\bottomrule
\end{tabular}
\end{center}

\noindent \textbf{Fabric} exposes Azure AI services natively in notebooks and pipelines.
\textbf{Snowflake} offers Document AI and Cortex functions---extraction as SQL.
\textbf{MongoDB} stores the extracted JSON naturally; Atlas Vector Search serves the retrieval
side of the same estate.

\section{Identity, Networking, and Key Management}
\index{IAM!equivalence}\index{private networking!equivalence}%
Chapter~21's security substrate---identity, private networking, key management---is identical in
concept on every cloud; only the acronyms differ.

\begin{center}
\footnotesize
\begin{tabular}{@{}p{2.4cm}p{2.9cm}p{2.9cm}p{2.9cm}p{2.9cm}@{}}
\toprule
\textbf{Capability} & \textbf{AWS} & \textbf{Azure} & \textbf{GCP} & \textbf{Databricks} \\
\midrule
Identity and access & IAM roles/policies & Entra ID + Azure RBAC & Cloud IAM & UC grants + SSO/SCIM \\
Org-wide guardrails & SCPs (Organizations) & Azure Policy & Organization Policy & Account-level policies + workspaces \\
Permissions ceiling & Permissions boundaries & Management groups + deny & IAM deny policies & Cluster policies + UC securables \\
Private service access & VPC endpoints / PrivateLink & Private Endpoints / Private Link & Private Service Connect & Private connectivity (control + data plane) \\
Key management & AWS KMS & Azure Key Vault & Cloud KMS & Cloud KMS (CMK on workspace storage) \\
Secrets management & Secrets Manager & Key Vault secrets & Secret Manager & Secret scopes (DB- or vault-backed) \\
\addlinespace
Free S3-class endpoint & Gateway endpoint & Service endpoints & Private Google Access & (via cloud provider endpoint) \\
Anti-exfiltration primitive & Endpoint policy + SCP + org ID & Private Link + policy & VPC Service Controls & Egress controls + storage credentials \\
\bottomrule
\end{tabular}
\end{center}

\noindent \textbf{Snowflake} adds account-level network policies and external-private-endpoint
rules; \textbf{MongoDB} Atlas uses private endpoints and IP access lists; \textbf{Fabric}
centralizes on Entra conditional access and tenant-level perimeter settings. The
pattern---identity plus network plus keys---is universal.

\section{Data Quality, Lineage, and Mesh Governance}
\index{data quality!equivalence}\index{lineage!equivalence}%
Chapter~22's governance layer---quality rules, lineage, mesh organization---converged faster here
than anywhere else in the stack.

\begin{center}
\footnotesize
\begin{tabular}{@{}p{2.4cm}p{2.9cm}p{2.9cm}p{2.9cm}p{2.9cm}@{}}
\toprule
\textbf{Capability} & \textbf{AWS} & \textbf{Azure} & \textbf{GCP} & \textbf{Databricks} \\
\midrule
Quality rules engine & Glue Data Quality (DQDL) & Fabric/Purview data quality & Dataplex data quality & DLT expectations (warn/drop/fail) \\
Cross-account sharing & Lake Formation + RAM & Purview + Fabric sharing & Analytics Hub / Dataplex & Delta Sharing \\
Tag-based governance & LF-Tags & Purview classifications & Policy tags & UC tags + governed tags \\
Lineage & SageMaker Lineage, OpenLineage & Purview lineage & Dataplex lineage & Automatic UC table/column lineage \\
Catalog of record & Glue Data Catalog & Microsoft Purview & Dataplex & Unity Catalog \\
\addlinespace
Open-source alternative & Great Expectations, DataHub & Same (cloud-agnostic) & Same & Same (+ MLflow for models) \\
Quality gate pattern & In-job evaluation + SNS & Pipeline gates & Dataflow/Dataplex gates & Expectation fail $\rightarrow$ pipeline halt \\
\bottomrule
\end{tabular}
\end{center}

\noindent \textbf{Snowflake} implements quality as data metric functions and governance as the
Horizon Catalog. \textbf{MongoDB} enforces shape at write time with schema validation---the
closest operational-store analog to a DQ ruleset. \textbf{Fabric} ties quality and lineage into
OneLake's unified governance story.

\section{Sensitive-Data Discovery and Erasure}
\index{PII discovery!equivalence}\index{erasure!equivalence}%
Chapter~23's regulatory constants---discovery, de-identification, erasure---get the same
architectural answers on every cloud.

\begin{center}
\footnotesize
\begin{tabular}{@{}p{2.4cm}p{2.9cm}p{2.9cm}p{2.9cm}p{2.9cm}@{}}
\toprule
\textbf{Capability} & \textbf{AWS} & \textbf{Azure} & \textbf{GCP} & \textbf{Databricks} \\
\midrule
PII discovery at rest & Amazon Macie & Microsoft Purview scanning & Sensitive Data Protection (DLP) & UC discovery (classification) \\
Custom identifiers & Macie custom data identifiers & Custom classifiers & Custom infoTypes & Custom classifiers / patterns \\
Audit trail & CloudTrail (incl.\ data events) & Azure Monitor / audit logs & Cloud Audit Logs & \texttt{system.access.audit} \\
Config posture & AWS Config rules & Azure Policy & Security Command Center & Account compliance + policies \\
Crypto-shredding & Per-subject KMS keys, then delete key & Per-subject KEKs in Key Vault & Per-subject keys in Cloud KMS & Per-subject keys via cloud KMS \\
Quarantine pattern & Object Lock bucket + Lambda & Immutability + Functions & Bucket lock + functions & Quarantine schema + UC grants \\
\addlinespace
Erasure across copies & Enumerate zones/snapshots/streams & Same problem, same pattern & Same problem, same pattern & Delete + VACUUM within retention window \\
\bottomrule
\end{tabular}
\end{center}

\noindent \textbf{Snowflake} offers object-tagging-driven masking and supports erasure via row
deletion within retention windows---with Time Travel/Fail-safe explicitly in scope.
\textbf{MongoDB} supports client-side field-level encryption, which makes per-subject
crypto-shredding natural. \textbf{Fabric} centralizes discovery in Purview over OneLake. On
Databricks, erasure must account for Delta history: the row delete and the \texttt{VACUUM} that
destroys old versions are one operation from the auditor's perspective.

\section{DataOps: IaC, Transformation-as-Code, and CI/CD}
\index{DataOps!equivalence}\index{Terraform!equivalence}\index{dbt!equivalence}%
Chapter~24 is the most portable week in the book: the tools are deliberately cloud-agnostic, so
the table maps providers rather than products.

\begin{center}
\footnotesize
\begin{tabular}{@{}p{2.4cm}p{2.9cm}p{2.9cm}p{2.9cm}p{2.9cm}@{}}
\toprule
\textbf{Capability} & \textbf{AWS} & \textbf{Azure} & \textbf{GCP} & \textbf{Databricks} \\
\midrule
IaC (multi-cloud) & Terraform AWS provider & Terraform azurerm & Terraform google & Terraform databricks provider \\
IaC (native) & CloudFormation / CDK & Bicep / ARM & Deployment Manager & Asset Bundles (\texttt{databricks.yml}) \\
Transformation as code & dbt-redshift & dbt-fabric / dbt-synapse & dbt-bigquery & dbt-databricks \\
CI/CD & GitHub Actions / CodePipeline & GitHub Actions / Azure DevOps & GitHub Actions / Cloud Build & GitHub Actions + bundle deploy \\
Remote state locking & S3 + DynamoDB & Azure Blob leases & GCS + state locking & (cloud backend of choice) \\
\addlinespace
CI authentication & OIDC to IAM role & OIDC to Entra workload identity & OIDC to Workload Identity & Service principal + OAuth M2M \\
Cost telemetry & CUR 2.0 & Cost Management & Billing export & \texttt{system.billing.usage} \\
Zero-downtime DDL & Expand-contract + view shims & Same pattern & Same pattern & Same pattern (+ Delta evolution) \\
\bottomrule
\end{tabular}
\end{center}

\noindent \textbf{Snowflake} (dbt-snowflake, Terraform provider) and \textbf{MongoDB} (Atlas
provider, schema-versioned migrations) plug into the same toolchain---this week's skills transfer
without translation. The lesson that holds on all seven platforms: \emph{the merge is the deploy,
and the pipeline is the only writer to production}.

\begin{tipbox}
Use this appendix as a study aid for the book's lockstep, cross-cloud approach: when you learn a
concept on one platform, immediately locate its equivalent here and predict how the other six
would express the same idea. That habit turns seven separate vocabularies into one mental model.
\end{tipbox}
